% --- Archon physics formalization source begin ---
% archon:physics
% archon:covers PhyXMiniProblems/problem_phyx_mini_0625.lean
% archon:source-report reports/phyx_mini/problem_phyx_mini_0625.source.json

\chapter{Physics problem phyx\_mini\_0625}
\label{ch:PhyXMiniProblems_problem_phyx_mini_0625}

\paragraph{Problem source.}
\#\# Physical scenario

Consider a monatomic solid with a weakly bound cubic lattice, with each atom connected to six neighbors, each bond having a binding energy of \textbackslash{}(3.4 \textbackslash{}times 10\textasciicircum{}\{-3\}\textbackslash{},\textbackslash{}mathrm\{eV\}\textbackslash{}). When this solid melts, its latent heat of fusion goes directly into breaking the bonds between the atoms.[Hint: Show that in a simple cubic lattice (figure), there are *three* times as many bonds as there are atoms, when the number of atoms is large.]

\#\# Question

Estimate the latent heat of fusion for this solid, in \textbackslash{}(\textbackslash{}mathrm\{J/mol\}\textbackslash{}).

\#\# Answer choices

- A: A: 960J/mol

- B: B: 970J/mol

- C: C: 990J/mol

- D: D: 980J/mol

\#\# Figure caption (auxiliary; use the image as primary evidence)

The image depicts a 3D lattice structure composed of uniformly spaced yellow spheres connected by black lines. The spheres are arranged in a grid-like pattern forming layers, and the lines represent bonds between adjacent spheres. There are no additional elements, text, or objects present in the image. The structure resembles a crystal lattice model commonly used to represent atomic or molecular arrangements in solid-state physics or chemistry.

\#\# Dataset metadata

Quantum Phenomena; Physical Model Grounding Reasoning, Multi-Formula Reasoning

\paragraph{Recorded answer/context.}
D: D: 980J/mol

\paragraph{Figure/image path.}
/root/proposal\_for\_physic/science-mango/phyx\_mini\_run/phyx\_data/test\_image/625.png

\paragraph{Formalization target.}
create a compiling Lean file with sorry bodies at `PhyXMiniProblems/problem\_phyx\_mini\_0625.lean`. The Lean declarations must preserve the physical quantities, dimensions or dimensional roles, figure labels, governing-law hypotheses, and final relation expressed by this problem.
Use Mathlib/Physlib names found through LeanExplore where available. If a physics API is missing, introduce faithful local abstractions rather than scalar placeholder aliases.

\begin{theorem}[Physics formalization target]
\label{thm:physics:phyx_mini_0625:target}
\lean{PhyXMiniProblems.ProblemPhyXMini0625.problem_phyx_mini_0625}
\uses{def:physics:phyx-mini-0625:phyxminiproblems-problemphyxmini0625-molarenergyscale, def:physics:phyx-mini-0625:phyxminiproblems-problemphyxmini0625-monatomiccubicsolidsetup, def:physics:phyx-mini-0625:phyxminiproblems-problemphyxmini0625-matchesweaklyboundsimplecubicscenario, def:physics:phyx-mini-0625:phyxminiproblems-problemphyxmini0625-matchessuppliedcubiclatticefigure, def:physics:phyx-mini-0625:phyxminiproblems-problemphyxmini0625-hasstatedbondandmolecalibrations, def:physics:phyx-mini-0625:phyxminiproblems-problemphyxmini0625-hasphysicalcubicsolidparameters, def:physics:phyx-mini-0625:phyxminiproblems-problemphyxmini0625-satisfiesfusionbondbreakinglaw, def:physics:phyx-mini-0625:phyxminiproblems-problemphyxmini0625-isuniqueclosestdisplayedchoice}
The assigned autoformalize agent should translate this physics problem into Lean declarations in the covered file, with theorem and lemma proof bodies written as `by sorry`.
\end{theorem}
\begin{proof}
This is an autoformalization task, not a proof task. Produce statements that can later be proved by the physics prover without weakening the physical model.
\end{proof}

% --- Archon declaration topology begin ---
\section{Declaration topology}
\begin{definition}[Molar Energy Scale]
  \label{def:physics:phyx-mini-0625:phyxminiproblems-problemphyxmini0625-molarenergyscale}
  \lean{PhyXMiniProblems.ProblemPhyXMini0625.MolarEnergyScale}
  An abstract carrier for molar energies, together with its scalar readout in joules per mole. A molar energy is therefore not identified with ℝ
\end{definition}
\begin{proof}
  This is the declared model definition of Molar Energy Scale.
\end{proof}
\begin{definition}[energy In Joules]
  \label{def:physics:phyx-mini-0625:phyxminiproblems-problemphyxmini0625-energyinjoules}
  \lean{PhyXMiniProblems.ProblemPhyXMini0625.energyInJoules}
  Read a dimensionful energy in coherent-SI joules
\end{definition}
\begin{proof}
  This is the declared model definition of energy In Joules.
\end{proof}
\begin{definition}[electron Volt In Joules]
  \label{def:physics:phyx-mini-0625:phyxminiproblems-problemphyxmini0625-electronvoltinjoules}
  \lean{PhyXMiniProblems.ProblemPhyXMini0625.electronVoltInJoules}
  \uses{def:physics:phyx-mini-0625:phyxminiproblems-problemphyxmini0625-energyinjoules}
  The coherent-SI joule value of Physlib's exact electron volt
\end{definition}
\begin{proof}
  This is the declared model definition of electron Volt In Joules.
\end{proof}
\begin{definition}[energy In Electron Volts]
  \label{def:physics:phyx-mini-0625:phyxminiproblems-problemphyxmini0625-energyinelectronvolts}
  \lean{PhyXMiniProblems.ProblemPhyXMini0625.energyInElectronVolts}
  \uses{def:physics:phyx-mini-0625:phyxminiproblems-problemphyxmini0625-energyinjoules, def:physics:phyx-mini-0625:phyxminiproblems-problemphyxmini0625-electronvoltinjoules}
  Read a physical energy in electron volts
\end{definition}
\begin{proof}
  This is the declared model definition of energy In Electron Volts.
\end{proof}
\begin{definition}[Solid Composition]
  \label{def:physics:phyx-mini-0625:phyxminiproblems-problemphyxmini0625-solidcomposition}
  \lean{PhyXMiniProblems.ProblemPhyXMini0625.SolidComposition}
  Physical composition of the solid
\end{definition}
\begin{proof}
  This is the declared model definition of Solid Composition.
\end{proof}
\begin{definition}[Lattice Kind]
  \label{def:physics:phyx-mini-0625:phyxminiproblems-problemphyxmini0625-latticekind}
  \lean{PhyXMiniProblems.ProblemPhyXMini0625.LatticeKind}
  Crystal-lattice class used by the physical model
\end{definition}
\begin{proof}
  This is the declared model definition of Lattice Kind.
\end{proof}
\begin{definition}[Binding Regime]
  \label{def:physics:phyx-mini-0625:phyxminiproblems-problemphyxmini0625-bindingregime}
  \lean{PhyXMiniProblems.ProblemPhyXMini0625.BindingRegime}
  Qualitative binding regime stated in the problem
\end{definition}
\begin{proof}
  This is the declared model definition of Binding Regime.
\end{proof}
\begin{definition}[Fusion Energy Destination]
  \label{def:physics:phyx-mini-0625:phyxminiproblems-problemphyxmini0625-fusionenergydestination}
  \lean{PhyXMiniProblems.ProblemPhyXMini0625.FusionEnergyDestination}
  Where the latent heat supplied during fusion is deposited
\end{definition}
\begin{proof}
  This is the declared model definition of Fusion Energy Destination.
\end{proof}
\begin{definition}[Figure Arrangement]
  \label{def:physics:phyx-mini-0625:phyxminiproblems-problemphyxmini0625-figurearrangement}
  \lean{PhyXMiniProblems.ProblemPhyXMini0625.FigureArrangement}
  Geometric arrangement visible in the supplied raster
\end{definition}
\begin{proof}
  This is the declared model definition of Figure Arrangement.
\end{proof}
\begin{definition}[Site Glyph]
  \label{def:physics:phyx-mini-0625:phyxminiproblems-problemphyxmini0625-siteglyph}
  \lean{PhyXMiniProblems.ProblemPhyXMini0625.SiteGlyph}
  Glyph used to depict an atom site
\end{definition}
\begin{proof}
  This is the declared model definition of Site Glyph.
\end{proof}
\begin{definition}[Figure Color]
  \label{def:physics:phyx-mini-0625:phyxminiproblems-problemphyxmini0625-figurecolor}
  \lean{PhyXMiniProblems.ProblemPhyXMini0625.FigureColor}
  Colors occurring in the raster
\end{definition}
\begin{proof}
  This is the declared model definition of Figure Color.
\end{proof}
\begin{definition}[Bond Stroke]
  \label{def:physics:phyx-mini-0625:phyxminiproblems-problemphyxmini0625-bondstroke}
  \lean{PhyXMiniProblems.ProblemPhyXMini0625.BondStroke}
  Stroke used to depict a bond
\end{definition}
\begin{proof}
  This is the declared model definition of Bond Stroke.
\end{proof}
\begin{definition}[Cubic Lattice Figure]
  \label{def:physics:phyx-mini-0625:phyxminiproblems-problemphyxmini0625-cubiclatticefigure}
  \lean{PhyXMiniProblems.ProblemPhyXMini0625.CubicLatticeFigure}
  \uses{def:physics:phyx-mini-0625:phyxminiproblems-problemphyxmini0625-figurearrangement, def:physics:phyx-mini-0625:phyxminiproblems-problemphyxmini0625-siteglyph, def:physics:phyx-mini-0625:phyxminiproblems-problemphyxmini0625-figurecolor, def:physics:phyx-mini-0625:phyxminiproblems-problemphyxmini0625-bondstroke}
  Literal qualitative data in the primary raster. The Boolean fields record visible presentation facts, not numerical bond counts or thermodynamic conclusions
\end{definition}
\begin{proof}
  This is the declared model definition of Cubic Lattice Figure.
\end{proof}
\begin{definition}[Monatomic Cubic Solid Setup]
  \label{def:physics:phyx-mini-0625:phyxminiproblems-problemphyxmini0625-monatomiccubicsolidsetup}
  \lean{PhyXMiniProblems.ProblemPhyXMini0625.MonatomicCubicSolidSetup}
  \uses{def:physics:phyx-mini-0625:phyxminiproblems-problemphyxmini0625-molarenergyscale, def:physics:phyx-mini-0625:phyxminiproblems-problemphyxmini0625-solidcomposition, def:physics:phyx-mini-0625:phyxminiproblems-problemphyxmini0625-latticekind, def:physics:phyx-mini-0625:phyxminiproblems-problemphyxmini0625-bindingregime, def:physics:phyx-mini-0625:phyxminiproblems-problemphyxmini0625-fusionenergydestination, def:physics:phyx-mini-0625:phyxminiproblems-problemphyxmini0625-cubiclatticefigure}
  Independent physical data for the solid. Atom indexes a representative finite bulk sample, and latticeBondGraph is its loopless undirected bond network. The latent heat is an independent quantity, not a definition made from an answer choice
\end{definition}
\begin{proof}
  This is the declared model definition of Monatomic Cubic Solid Setup.
\end{proof}
\begin{definition}[Matches Weakly Bound Simple Cubic Scenario]
  \label{def:physics:phyx-mini-0625:phyxminiproblems-problemphyxmini0625-matchesweaklyboundsimplecubicscenario}
  \lean{PhyXMiniProblems.ProblemPhyXMini0625.MatchesWeaklyBoundSimpleCubicScenario}
  \uses{def:physics:phyx-mini-0625:phyxminiproblems-problemphyxmini0625-molarenergyscale, def:physics:phyx-mini-0625:phyxminiproblems-problemphyxmini0625-monatomiccubicsolidsetup}
  The prose-level physical scenario. Six-regularity is the bulk nearest-neighbour idealization. Since SimpleGraph is loopless and undirected, one bond joins two distinct sites and is not double-oriented
\end{definition}
\begin{proof}
  This is the declared model definition of Matches Weakly Bound Simple Cubic Scenario.
\end{proof}
\begin{definition}[Matches Supplied Cubic Lattice Figure]
  \label{def:physics:phyx-mini-0625:phyxminiproblems-problemphyxmini0625-matchessuppliedcubiclatticefigure}
  \lean{PhyXMiniProblems.ProblemPhyXMini0625.MatchesSuppliedCubicLatticeFigure}
  \uses{def:physics:phyx-mini-0625:phyxminiproblems-problemphyxmini0625-molarenergyscale, def:physics:phyx-mini-0625:phyxminiproblems-problemphyxmini0625-monatomiccubicsolidsetup}
  Facts read from the primary raster. In particular, the image has no text, numeric answer, or thermodynamic annotation
\end{definition}
\begin{proof}
  This is the declared model definition of Matches Supplied Cubic Lattice Figure.
\end{proof}
\begin{definition}[Has Stated Bond And Mole Calibrations]
  \label{def:physics:phyx-mini-0625:phyxminiproblems-problemphyxmini0625-hasstatedbondandmolecalibrations}
  \lean{PhyXMiniProblems.ProblemPhyXMini0625.HasStatedBondAndMoleCalibrations}
  \uses{def:physics:phyx-mini-0625:phyxminiproblems-problemphyxmini0625-molarenergyscale, def:physics:phyx-mini-0625:phyxminiproblems-problemphyxmini0625-energyinelectronvolts, def:physics:phyx-mini-0625:phyxminiproblems-problemphyxmini0625-monatomiccubicsolidsetup}
  Stated numerical inputs at explicit unit-readout boundaries. The bond energy is 3.4 * 10⁻³ eV, and the second field is the exact number of atoms per mole. Neither field mentions the requested latent heat
\end{definition}
\begin{proof}
  This is the declared model definition of Has Stated Bond And Mole Calibrations.
\end{proof}
\begin{definition}[Has Physical Cubic Solid Parameters]
  \label{def:physics:phyx-mini-0625:phyxminiproblems-problemphyxmini0625-hasphysicalcubicsolidparameters}
  \lean{PhyXMiniProblems.ProblemPhyXMini0625.HasPhysicalCubicSolidParameters}
  \uses{def:physics:phyx-mini-0625:phyxminiproblems-problemphyxmini0625-molarenergyscale, def:physics:phyx-mini-0625:phyxminiproblems-problemphyxmini0625-energyinjoules, def:physics:phyx-mini-0625:phyxminiproblems-problemphyxmini0625-monatomiccubicsolidsetup}
  Positivity and non-vacuity conditions for a physical bulk sample
\end{definition}
\begin{proof}
  This is the declared model definition of Has Physical Cubic Solid Parameters.
\end{proof}
\begin{definition}[Satisfies Fusion Bond Breaking Law]
  \label{def:physics:phyx-mini-0625:phyxminiproblems-problemphyxmini0625-satisfiesfusionbondbreakinglaw}
  \lean{PhyXMiniProblems.ProblemPhyXMini0625.SatisfiesFusionBondBreakingLaw}
  \uses{def:physics:phyx-mini-0625:phyxminiproblems-problemphyxmini0625-molarenergyscale, def:physics:phyx-mini-0625:phyxminiproblems-problemphyxmini0625-energyinjoules, def:physics:phyx-mini-0625:phyxminiproblems-problemphyxmini0625-monatomiccubicsolidsetup}
  General bond-breaking energy accounting for one mole. Multiplication by the number of atoms in the representative bulk sample avoids assuming the desired three-bonds-per-atom conclusion; that relation must be derived from six-regularity by the handshaking lemma
\end{definition}
\begin{proof}
  This is the declared model definition of Satisfies Fusion Bond Breaking Law.
\end{proof}
\begin{lemma}[bulk Bond Count eq three mul atom Count]
  \label{lem:physics:phyx-mini-0625:phyxminiproblems-problemphyxmini0625-bulkbondcount-eq-three-mul-atomcount}
  \lean{PhyXMiniProblems.ProblemPhyXMini0625.bulkBondCount_eq_three_mul_atomCount}
  The hint's bond-count conclusion: a finite six-regular bulk model has three times as many undirected bonds as atom sites
\end{lemma}
\begin{proof}
  The conclusion follows from the governing relations and figure readouts named in the statement of bulk Bond Count eq three mul atom Count.
\end{proof}
\begin{lemma}[latent Heat eq three bonds per atom]
  \label{lem:physics:phyx-mini-0625:phyxminiproblems-problemphyxmini0625-latentheat-eq-three-bonds-per-atom}
  \lean{PhyXMiniProblems.ProblemPhyXMini0625.latentHeat_eq_three_bonds_per_atom}
  \uses{def:physics:phyx-mini-0625:phyxminiproblems-problemphyxmini0625-molarenergyscale, def:physics:phyx-mini-0625:phyxminiproblems-problemphyxmini0625-energyinjoules, def:physics:phyx-mini-0625:phyxminiproblems-problemphyxmini0625-monatomiccubicsolidsetup, def:physics:phyx-mini-0625:phyxminiproblems-problemphyxmini0625-matchesweaklyboundsimplecubicscenario, def:physics:phyx-mini-0625:phyxminiproblems-problemphyxmini0625-hasphysicalcubicsolidparameters, def:physics:phyx-mini-0625:phyxminiproblems-problemphyxmini0625-satisfiesfusionbondbreakinglaw}
  The energy balance and derived bond count express molar latent heat as three bond energies per atom times the number of atoms in one mole
\end{lemma}
\begin{proof}
  The conclusion follows from the governing relations and figure readouts named in the statement of latent Heat eq three bonds per atom.
\end{proof}
\begin{definition}[Answer Choice]
  \label{def:physics:phyx-mini-0625:phyxminiproblems-problemphyxmini0625-answerchoice}
  \lean{PhyXMiniProblems.ProblemPhyXMini0625.AnswerChoice}
  Labels of the four displayed answer choices
\end{definition}
\begin{proof}
  This is the declared model definition of Answer Choice.
\end{proof}
\begin{definition}[displayed Latent Heat In Joules Per Mole]
  \label{def:physics:phyx-mini-0625:phyxminiproblems-problemphyxmini0625-displayedlatentheatinjoulespermole}
  \lean{PhyXMiniProblems.ProblemPhyXMini0625.displayedLatentHeatInJoulesPerMole}
  \uses{def:physics:phyx-mini-0625:phyxminiproblems-problemphyxmini0625-answerchoice}
  Displayed latent heats of fusion, read in joules per mole
\end{definition}
\begin{proof}
  This is the declared model definition of displayed Latent Heat In Joules Per Mole.
\end{proof}
\begin{definition}[recorded Answer]
  \label{def:physics:phyx-mini-0625:phyxminiproblems-problemphyxmini0625-recordedanswer}
  \lean{PhyXMiniProblems.ProblemPhyXMini0625.recordedAnswer}
  \uses{def:physics:phyx-mini-0625:phyxminiproblems-problemphyxmini0625-answerchoice}
  The source dataset records choice D; this metadata is not a premise
\end{definition}
\begin{proof}
  This is the declared model definition of recorded Answer.
\end{proof}
\begin{definition}[Is Unique Closest Displayed Choice]
  \label{def:physics:phyx-mini-0625:phyxminiproblems-problemphyxmini0625-isuniqueclosestdisplayedchoice}
  \lean{PhyXMiniProblems.ProblemPhyXMini0625.IsUniqueClosestDisplayedChoice}
  \uses{def:physics:phyx-mini-0625:phyxminiproblems-problemphyxmini0625-molarenergyscale, def:physics:phyx-mini-0625:phyxminiproblems-problemphyxmini0625-answerchoice, def:physics:phyx-mini-0625:phyxminiproblems-problemphyxmini0625-displayedlatentheatinjoulespermole}
  A displayed choice is uniquely closest to the physical molar-latent-heat readout. Strict comparison is required only against different labels
\end{definition}
\begin{proof}
  This is the declared model definition of Is Unique Closest Displayed Choice.
\end{proof}
% --- Archon declaration topology end ---
% --- Archon physics formalization source end ---
